% SGA 6，第0讲附录 RRR；PDF 30--57（止于第二章第3节之前）。
\SourcePage{30}
\SGAchapter{0，附录 RRR}{层的类与黎曼—罗赫定理}{Classes de faisceaux et théorème de Riemann--Roch}
\begin{center}A. Grothendieck\end{center}
\footnote{这是本研讨班引言中所引报告的逐字重印。以（*）、（**）标出的脚注于1967年10月补入。}
原文目录如下（\Term{lie-bracket}{括号}内为报告原页码）：第一章，$\lambda$-环（\Term{form}{形式}预备知识）：1. 定义（1）；2. 例子（4）；3. \Term{graded-ring}{分次环}所定义的$\lambda$-环（8）；4. 运算$\lambda^p(N,x)$（13）。第二章，\Term{coherent-algebraic-sheaf}{凝聚代数层}的类与\Term{chern-class}{陈类}：1. 周理论（18）；2. \Term{coherent-algebraic-sheaf}{凝聚代数层}的\Term{chern-class}{陈类}的定义（22）；3. 关于$K(X)$的\Term{functorial}{函子性}的\Term{generic}{一般}事项（27）；4. 若干技术结果（34）；5. \Term{chern-class}{陈类}的层论定义。对\Term{embedding}{嵌入}\Term{morphism}{态射}研究的应用（43）；6. 黎曼—罗赫\Term{theorem}{定理}（49）。另附：黎曼—罗赫\Term{theorem}{定理}的证明（1957年11月1日）。

\SourcePage{31}
\SGAchapter{RRR I}{$\lambda$-环（形式预备知识）}{$\lambda$-Anneaux (Préliminaires formels)}
\SGAsection{1}{定义}
所谓\Term{lambda-ring}{$\lambda$-环}\footnote{本研讨班中，我们更常称之为\Term{pre-lambda-ring}{预$\lambda$-环}（第V讲2.1），而将\Term{lambda-ring}{$\lambda$-环}这一名称留给满足原报告第3页附加条件的环（参见第V讲2.4）。}，是指一个带\Term{unit}{单位}的\Term{commutative-ring}{交换环}$K$，配备一族映射
\begin{equation}\lambda^i:K\longrightarrow K\tag{1.1}\end{equation}
（$i$为整数且$i\geq0$），满足下列条件：
\begin{equation}\lambda^0x=1,\quad\lambda^1x=x,\quad
\lambda^n(x+y)=\sum_{i=0}^n\lambda^ix\,\lambda^{n-i}y.\tag{1.2}\end{equation}
对每个$x\in K$，置
\begin{equation}\lambda_t(x)=\sum_{n\geq0}\lambda^n(x)t^n\in K[[t]].\tag{1.3}\end{equation}
上述关系于是表明，$x\mapsto\lambda_t(x)$是从\Term{additive-group}{加法群}$K$到$K$中\Term{augmentation}{增广}为1的\Term{formal-power-series}{形式幂级数}的\Term{multiplicative-group}{乘法群}$1+K[[t]]^+$的\Term{homomorphism}{同态}，且为\Term{canonical}{典范}\Term{homomorphism}{同态}$1+\sum_{i>0}x_it^i\mapsto x_1$的\Term{lifting}{提升}。
因此，环$\mathbb Z$上存在唯一与其\Term{ring-structure}{环结构}相容并满足下式的$\lambda$-结构：
\SourcePage{32}
\begin{equation}\lambda_t(1)=1+t.\tag{1.4}\end{equation}
于是有
\begin{equation}\lambda_t(n)=(1+t)^n,\qquad\text{从而}\quad\lambda^i(n)=\binom ni.\tag{1.5}\end{equation}
\Term{lambda-ring}{$\lambda$-环}的\Term{lambda-homomorphism}{$\lambda$-同态}这一概念是清楚的；所谓\Term{augmented-lambda-ring}{增广$\lambda$-环}，是指配备一个到$\mathbb Z$的\Term{lambda-homomorphism}{$\lambda$-同态}的\Term{lambda-ring}{$\lambda$-环}。

为定义\Term{special-lambda-ring}{特殊$\lambda$-环}，下面的考虑是有用的。设$k$为带\Term{unit}{单位}的\Term{commutative-ring}{交换环}，$K=1+k[[t]]^+$为系数在$k$中、\Term{augmentation}{增广}为1的\Term{formal-power-series}{形式幂级数}群；我们要在其上引入一个\Term{lambda-ring}{$\lambda$-环}运算，使其加法结构就是\Term{formal-power-series}{形式幂级数}的\Term{multiplication}{乘法}。这个环的\Term{multiplication}{乘法}记为$f\circ g$。它由以下要求完全确定：$f\circ g$的系数$c_n=c_n(f\circ g)$由分别为$f,g$的系数$a_i=c_i(f)$、$b_i=c_i(g)$的整系数\Term{universal-polynomial}{泛多项式}给出：
\begin{equation}c_n=P_n(a_1,\ldots,a_n;b_1,\ldots,b_n);\tag{1.6}\end{equation}
它对“加法”$fg$满足\Term{distributivity}{分配律}，并且对一次因子有
\begin{equation}(1+at)\circ(1+bt)=1+abt.\tag{1.7}\end{equation}
$P_n$可用\Term{symmetric-function}{对称函数}的计算求出；$P_n$关于$a_i$和关于$b_i$都为权$n$的\Term{isobaric-polynomial}{等权多项式}（$a_i,b_i$的权为$i$），而且关于$a=(a_i)$与$b=(b_i)$\Term{symmetric}{对称}。同样，$\lambda^i(f)$由以下要求唯一确定：
\SourcePage{33}
$\lambda^i(f)$的系数$c_n(\lambda^i(f))=c_n$可由$f$的系数$c_n(f)=a_n$的整系数\Term{universal-polynomial}{泛多项式}算出：
\begin{equation}c_n=P_{i,n}(a_1,\ldots,a_{in}),\tag{1.8}\end{equation}
关系(1.2)成立，而且对$i>1$有
\begin{equation}\lambda^i(1+at)=1\tag{1.9}\end{equation}
（$1$为环$K=1+k[[t]]^+$的零元）。$P_{i,n}$仍可用\Term{symmetric-function}{对称函数}计算求出，并可看出它是关于$a_i$的权$in$的\Term{isobaric-polynomial}{等权多项式}。配备上述运算后，$K$成为一个\Term{lambda-ring}{$\lambda$-环}，其零元为1，\Term{identity-element}{单位元}为$1+t$。还应注意关系
\begin{equation}(1+at)\circ f=f(at).\tag{1.7 bis}\end{equation}
现在设$K$为任意\Term{lambda-ring}{$\lambda$-环}。若加法\Term{homomorphism}{同态}$\lambda_t:K\to1+K[[t]]^+$为\Term{lambda-homomorphism}{$\lambda$-同态}，就称$K$为\Term{special-lambda-ring}{特殊的}。这就是说，有公式
\begin{equation}
\lambda_t(1)=1+t,\quad\lambda_t(xy)=\lambda_t(x)\circ\lambda_t(y),\quad
\lambda_t(\lambda^i(x))=\lambda^i(\lambda_t(x)).\tag{1.10}
\end{equation}
显式写出即为
\begin{equation}\left\{\begin{aligned}
\lambda^i(1)&=0&&\text{若 }i\geq1,\\
\lambda^n(xy)&=P_n(\lambda^1x,\ldots,\lambda^nx;\lambda^1y,\ldots,\lambda^ny),\\
\lambda^n(\lambda^ix)&=P_{i,n}(\lambda^1x,\ldots,\lambda^nx).
\end{aligned}\right.\tag{1.11}\end{equation}
其中$P_n$与$P_{i,n}$为上面所述的\Term{universal-polynomial}{泛多项式}。特别地，由此可得公式(1.5)。（可以证明，从\Term{commutative-ring}{交换环}$k$构造的\Term{lambda-ring}{$\lambda$-环}$1+k[[t]]^+$是特殊的；\SourcePage{34} 但我们无需使用这一事实。）

\SGAsection{2}{例子}
1）设$G$为群，$k$为\Term{commutative-ring}{交换环}；考虑群$\underline K(G)$，它是以$G$在\Term{finite-type-module}{有限型$k$-模}中的\Term{representation}{表示}的类为基的\Term{free-group}{自由群}，模去形如$(\rho\oplus\rho')-(\rho)-(\rho')$的元素所\Term{generated-subgroup}{生成的子群}所得的\Term{quotient-group}{商群}（$\oplus$表示\Term{direct-sum}{直和}）。\Term{tensor-product}{张量积}与\Term{exterior-power}{外幂}在$\underline K(G)$上定义一个\Term{lambda-ring}{$\lambda$-环}结构。还可以再模去形如$(\rho)-(\rho')-(\rho'')$的元素，其中$\rho$是由$\rho''$\Term{extension}{扩张}$\rho'$得到的表示，且作为$k$-模的\Term{extension}{扩张}是平凡的。这样得到$G$的表示的\Term{reduced-class}{约化类}群$\underline K_r(G)$；可以看出，\Term{lambda-ring}{$\lambda$-环}运算\Term{descent}{下降}到此\Term{quotient-group}{商群}上。

当$k$为\Term{algebraically-closed-field}{代数闭域}且\Term{characteristic}{特征}为0时\footnote{$k$为\Term{algebraically-closed-field}{代数闭域}这一要求是多余的；参见第VI讲3.3中更\Term{generic}{一般}的结果，它包括本报告将使用的所有结果。}，可证明$\underline K(G)$，从而更有$\underline K_r(G)$，为\Term{special-lambda-ring}{特殊$\lambda$-环}；\Term{characteristic}{特征}不为0时不再如此，但若$k$为\Term{algebraically-closed-field}{代数闭域}（*），$\underline K_r(G)$仍是\Term{special-lambda-ring}{特殊$\lambda$-环}。为证明这一点，可归结到$G$为乘积$\operatorname{GL}(m,k)\times\operatorname{GL}(n,k)$，且只考虑$G$的\Term{rational-representation}{有理表示}的情形；我们的断言由下面例3中对$\underline K_r(G)$的显式确定得出。最后，当$k$的\Term{characteristic}{特征}为0时，$G$的\Term{rational-representation}{有理表示}的\Term{complete-reducibility}{完全可约性}表明$\underline K(G)=\underline K_r(G)$。

\SourcePage{35}
前例的变体：$G$为定义在$k$上的\Term{algebraic-group}{代数群}，且只考虑$G$的定义在$k$上的\Term{rational-representation}{有理表示}；也可以考虑\Term{topological-group}{拓扑群}或\Term{lie-group}{李群}及其\Term{continuous-representation}{连续表示}，等等。\footnote{还可以构造其他\Term{lambda-ring}{$\lambda$-环}的例子，例如\Term{quadratic-form}{二次型}的\Term{equivalence-class}{等价类}的\Term{set}{集合}、\Term{lie-algebra}{李代数}的表示类的\Term{set}{集合}，等等……我们保证，处理黎曼—罗赫\Term{theorem}{定理}并不需要所有这些例子！}

2）设$X$为定义在域$k$上的\Term{algebraic-variety}{代数簇}，无论\Term{reducible}{可约}与否。可以考虑以$X$上（定义在$k$上的）\Term{vector-bundle}{向量丛}的类为基的\Term{free-group}{自由群}，模去由形如$(E\oplus E')-(E)-(E')$的元素\Term{generated-subgroup}{生成的子群}所得的\Term{quotient-group}{商群}；\Term{tensor-product}{张量积}与\Term{exterior-power}{外幂}在其上定义\Term{lambda-ring}{$\lambda$-环}运算。若$X$是\Term{complete-variety}{完备的}，则由Krull--Remak--Schmidt\Term{theorem}{定理}，这是以$X$上的\Term{indecomposable-vector-bundle}{不可分解向量丛}为基的\Term{free-group}{自由群}；知道这个环及其基，\Term{equivalence}{等价}于知道$X$上的\Term{vector-bundle}{向量丛}的完全分类，以及这些丛上的基本\Term{tensor-operation}{张量运算}。若$k$为\Term{characteristic}{特征}0的\Term{algebraically-closed-field}{代数闭域}，则借助以下3）的内容，也就知道了丛的类上的所有其他\Term{tensor-operation}{张量运算}。还可以作更彻底的取商：模去由形如$(E)-(E')-(E'')$的元素生成的群，其中\Term{vector-bundle}{向量丛}$E$是由$E''$\Term{extension}{扩张}$E'$得到的；于是得到一个\Term{lambda-ring}{$\lambda$-环}，记为$\underline K(\mathscr S(X))$，其中$\mathscr S(X)$表示$X$上定义在$k$上的\Term{vector-bundle}{向量丛}构成的\Term{additive-category}{加性范畴}。用例1的结果可证明：若$k$为\Term{algebraically-closed-field}{代数闭域}\footnote{这一条件是多余的，参见上一脚注。}，则\Term{lambda-ring}{$\lambda$-环}$\underline K(\mathscr S(X))$是特殊的；\SourcePage{36} 当$k$的\Term{characteristic}{特征}为0时，这里定义的第一个\Term{lambda-ring}{$\lambda$-环}也如此，但\Term{generic}{一般}情形并非如此。

3）在若干重要情形中确定$\underline K(G)$与$\underline K_r(G)$。设域$k$为\Term{algebraically-closed-field}{代数闭域}，群$G$为\Term{affine-algebraic-group}{仿射代数群}、\Term{connected}{连通}且“\Term{globally-reductive}{整体约化}”（$G$的\Term{radical}{根基}\Term{isomorphism}{同构}于$(k^*)^s$），并且只考虑有理\Term{linear-representation}{线性表示}。\Term{characteristic}{特征}0时，有理\Term{linear-representation}{线性表示}的\Term{complete-reducibility}{完全可约性}表明，$\underline K(G)=\underline K_r(G)$是以\Term{irreducible}{不可约}\Term{rational-representation}{有理表示}的类为基的\Term{free-group}{自由群}。若$T$为$G$的一个\Term{maximal-torus}{极大环面}，则$\underline K(T)=\underline K_r(T)=\mathbb Z(\widehat T)$，即$T$的\Term{rational-character}{有理特征}群$\widehat T$的\Term{group-algebra}{群代数}；若$r$为$G$的\Term{rank}{秩}，则$\widehat T\simeq\mathbb Z^r$（这与\Term{characteristic}{特征}无关）。\Term{weyl-group}{Weyl群}$W$\Term{action}{作用}在$T$上，因而\Term{action}{作用}在$\underline K(T)$上，其\Term{action}{作用}方式只取决于$W$在\Term{free-group}{自由群}$\widehat T$上的\Term{action}{作用}。由$T$到$G$的\Term{embedding}{嵌入}\Term{derive}{导出}显然的\Term{lambda-ring}{$\lambda$-环}\Term{homomorphism}{同态}
\begin{equation}\underline K_r(G)\longrightarrow\underline K_r(T)^W=\underline K(T)^W,\tag{1.13}\end{equation}
其中$\underline K(T)^W$表示$\underline K(T)$中$W$的\Term{fixed-point}{不动点}\Term{set}{集合}。$G$的\Term{linear-representation}{线性表示}理论\footnote{参见Chevalley研讨班1956/58，\textit{Groupes de Lie algébriques}（巴黎高等师范学校）。}给出下述\Term{theorem}{定理}。
\begin{SGAresult}{\Term{theorem}{定理}}{1.1}\Term{homomorphism}{同态}(1.13)是\Term{isomorphism}{同构}。\end{SGAresult}
因此，一旦知道$G$的“类型”（由$W$在群$\mathbb Z^r$上的\Term{action}{作用}来刻画），$\underline K_r(G)$就与\Term{characteristic}{特征}无关，而且是一个\Term{special-lambda-ring}{特殊$\lambda$-环}。$\underline K_r(G)$有一个与$W$在$\widehat T$中的\Term{orbit}{轨道}对应的\Term{canonical-basis}{典范基}；另一个基可如下得到：
\SourcePage{37}
设$(\rho_i)_{1\leq i\leq r'}$为与$G/\operatorname{rad}G$的基本表示对应的$G$的表示类，$(\sigma_i)_{1\leq i\leq r''}$为$G$到$k^*$的有理\Term{homomorphism}{同态}群的一组基；于是$r'+r''=r$（$G$的\Term{rank}{秩}）。
\begin{SGAresult}{推论}{}
环$\underline K_r(G)$可等同于以$\rho_i,\sigma_j$为变量、系数在$\mathbb Q$中的\Term{rational-function-field}{有理函数域}的\Term{subring}{子环}
\[\mathbb Z[\rho_1,\ldots,\rho_{r'};\sigma_1,\ldots,\sigma_{r''};\sigma_1^{-1},\ldots,\sigma_{r''}^{-1}].\]
\end{SGAresult}
\footnote{为使这一陈述正确，必须假定$G$的\Term{derived-group}{导出群}$G'$为\Term{simply-connected}{单连通}。}
特别地，$\underline K_r(G)$没有\Term{zero-divisor}{零因子}；若$G=\prod_i\operatorname{GL}(n_i,k)$，则环$\underline K_r(G)$可如下显式写出：设$\rho_i$为从$\operatorname{GL}(n_i,k)$的恒等表示\Term{derive}{导出}的$G$的表示；考虑以$\lambda^j(\rho_i)$（$1\leq j\leq n_i$）为\Term{indeterminate}{不定元}、系数在$\mathbb Q$中的\Term{rational-function-field}{有理函数域}，则$\underline K_r(G)$可等同于由这些\Term{indeterminate}{不定元}及$\lambda^{n_i}(\rho_i)$的逆元生成的\Term{subring}{子环}。这证实了我们的断言：至少在\Term{characteristic}{特征}0时，运算$\lambda^i$可重构\Term{vector-bundle}{向量丛}上的全部\Term{tensor-operation}{张量运算}；在任意\Term{characteristic}{特征}的\Term{algebraically-closed-field}{代数闭域}上，只要取\Term{reduced-class}{约化类}，这也仍然成立。

\textbf{注记。}a）例1中的环$\underline K(G)$带有\Term{augmentation}{增广}，其\Term{augmentation}{增广}由表示的\Term{degree}{次数}定义；通过取商，这一\Term{augmentation}{增广}在$\underline K_r(G)$上定义一个\Term{augmentation}{增广}。同样，当$X$为$k$-\Term{irreducible}{不可约}时，例2中考虑的\Term{lambda-ring}{$\lambda$-环}带有\Term{augmentation}{增广}；在$\underline K(\mathscr S(X))$中，\Term{augmentation}{增广}由\Term{vector-bundle}{向量丛}的\Term{rank}{秩}，即其\Term{fiber}{纤维}的\Term{dimension}{维数}定义。
\SourcePage{38}
b）当域$k$不是\Term{algebraically-closed-field}{代数闭域}时，似乎仍有：\Term{characteristic}{特征}0时$\underline K(G)$为\Term{special-lambda-ring}{特殊$\lambda$-环}，而任意\Term{characteristic}{特征}时$\underline K_r(G)$为\Term{special-lambda-ring}{特殊$\lambda$-环}。

c）在非零\Term{characteristic}{特征}的情形，\Term{theorem}{定理}1.1由Chevalley惠予提供。

\SGAsection{3}{分次环所定义的$\lambda$-环}
设$A$为带\Term{unit}{单位}的交换\Term{graded-ring}{分次环}；为简单起见，假定$A$的\Term{degree}{次数}非负，且$A^0=\mathbb Z$。以$\widehat A$表示各$A^i$的\Term{direct-product}{直积}（这是带\Term{unit}{单位}的\Term{commutative-ring}{交换环}），以$\widehat A^+$表示显然的\Term{augmentation}{增广}$\widehat A\to A^0=\mathbb Z$的核。于是$1+\widehat A^+$为$\widehat A$的\Term{invertible-element}{可逆元}\Term{multiplicative-group}{乘法群}的\Term{subgroup}{子群}。考虑\Term{abelian-group}{阿贝尔群}
\begin{equation}\widetilde A=\mathbb Z\times(1+\widehat A^+),\tag{1.14}\end{equation}
其元素记为$[n,x]$，其中$n\in\mathbb Z$，$x=1+\sum_{i\geq1}x^i\in1+\widehat A^+$（$x^i\in A^i$）。这个群采用加法记号，故
\begin{equation}[n,x]+[n',x']=[n+n',xx'].\tag{1.15}\end{equation}
它的零元为$[0,1]$。我们要在其上引入与其加法结构相容的\Term{augmented-lambda-ring}{增广$\lambda$-环}结构，\Term{augmentation}{增广}为$\epsilon:[n,x]\mapsto n$。乘积$[m,x][n,y]$由下列条件完全刻画：它写成$[mn,x^ny^m(x*y)]$，其中$x*y\in1+\widehat A^+$，其\Term{set-component}{分量}为$x^i,y^i$的整系数\Term{universal-polynomial}{泛多项式}：
\begin{equation}(x*y)^i=Q_i(x^1,\ldots,x^i;y^1,\ldots,y^i)\qquad(i\geq1),\tag{1.16}\end{equation}
\SourcePage{39}
\Term{polynomial}{多项式}$Q_i$为权$i$的\Term{isobaric-polynomial}{等权多项式}。此外，$x*y$须对“加法”$xy$满足\Term{distributivity}{分配律}，并且对于“一次”因子化为
\begin{equation}[1,1+x^1][1,1+y^1]=[1,1+x^1+y^1].\tag{1.17}\end{equation}
\Term{polynomial}{多项式}$Q_i$可由初等\Term{symmetric-function}{对称函数}的计算求出（这就是利用因子的\Term{chern-class}{陈类}$x^i,y^i$及各自的\Term{dimension}{维数}$m,n$，计算\Term{vector-bundle}{向量丛}\Term{tensor-product}{张量积}的\Term{chern-class}{陈类}）。由(1.17)得
\begin{equation}(1+x^1)*(1+y^1)=\frac{1+x^1+y^1}{(1+x^1)(1+y^1)}.\tag{1.17 bis}\end{equation}
将$x$“\Term{form}{形式}地”\Term{resolution}{分解}为乘积$\prod_{i=1}^n(1+\alpha_i)$，将$y$\Term{resolution}{分解}为$\prod_{i=1}^m(1+\beta_i)$，其中$\alpha_i,\beta_i$的\Term{degree}{次数}为1，就可逐项计算任意乘积$x*y$。

这样，$\widetilde A$成为一个带\Term{augmentation}{增广}的\Term{commutative-ring}{交换环}，\Term{identity-element}{单位元}为$[1,1]$；也可将它理解为在环$1+\widehat A^+$中添加\Term{identity-element}{单位元}所得的环。此外，在$\widetilde A$上定义\Term{filtration}{滤过}如下：置$\widetilde A^0=\widetilde A$，以$\widetilde A^n$表示所有$[0,x]$构成的\Term{set}{集合}，其中当$1\leq i<n$时$x^i=0$。这个\Term{filtration}{滤过}与\Term{multiplication}{乘法}相容，因为
\begin{equation}
\left(1+\sum_{i\geq m}x^i\right)*\left(1+\sum_{j\geq n}y^j\right)
=1-\frac{(m+n-1)!}{(m-1)!(n-1)!}x^my^n+\cdots,\tag{1.18}
\end{equation}
未写出的各项的\Term{degree}{次数}大于$m+n$。此式表明，与$\widetilde A$相伴的\Term{associated-graded-ring}{分次环}$G(\widetilde A)$在加法意义上可等同于$A$，但在\Term{multiplication}{乘法}意义上不能如此；我们只能说，以下定义的$A$到$G(\widetilde A)$的映射$\eta'$：
\SourcePage{40}
\begin{equation}\eta'(x^i)=(-1)^{i-1}(i-1)!\,x^i\qquad(x^i\in A^i)\tag{1.19}\end{equation}
是环\Term{homomorphism}{同态}。

在$\widetilde A$中引入满足\Term{lambda-ring}{$\lambda$-环}公理（第1节，公式(1.2)）的运算$\lambda^i$；它们由以下条件唯一确定：$\lambda^i[0,x]$的\Term{set-component}{分量}由$x$的\Term{set-component}{分量}$x^i$通过整系数\Term{universal-polynomial}{泛多项式}给出：
\begin{equation}\lambda^i[0,x]=[0,\lambda^ix],\qquad
(\lambda^ix)^{(n)}=Q_{i,n}(x^1,\ldots,x^n)\tag{1.20}\end{equation}
（$Q_{i,n}$为权$n$的\Term{isobaric-polynomial}{等权多项式}），并且
\begin{equation}\lambda^i[1,1+x^1]=0\qquad\text{对 }i>1.\tag{1.21}\end{equation}
\Term{polynomial}{多项式}$Q_{i,n}$由初等\Term{symmetric-function}{对称函数}的计算求出，即由\Term{vector-bundle}{向量丛}的\Term{chern-class}{陈类}和\Term{rank}{秩}计算其\Term{exterior-power}{外幂}的\Term{chern-class}{陈类}。于是可看出，$\widetilde A$为\Term{augmentation}{增广}的\Term{special-lambda-ring}{特殊$\lambda$-环}，且各$\widetilde A^i$在运算$\lambda^j$下\Term{stable}{稳定}。

\Term{augmented-lambda-ring}{增广$\lambda$-环}$\widetilde A$可视为$A$的一个\Term{functor}{函子}；因此$\widehat A$的\Term{principal-automorphism}{主自同构}
\[x=\sum_{i\geq0}x^i\longmapsto x^\vee=\sum_{i\geq0}(-1)^ix^i\]
定义$\widetilde A$的一个\Term{automorphism}{自同构}，仍用\Term{signature}{符号}$\vee$表示。
\begin{SGAresult}{命题}{1.2}
设$x\in\widetilde A$形如$[n,1+x^1+\cdots+x^n]$（故$i>n$时$x^i=0$）。则$i>n$时$\lambda^i(x)=0$，并有$\lambda^n(x)=[1,1+x^1]$；最后有
\begin{equation}\sum_{i=0}^n(-1)^i\lambda^i(x)=[0,1-(n-1)!x^n+\cdots].\tag{1.22}\end{equation}
\end{SGAresult}
由此命题以及各$\widetilde A^i$在运算$\lambda^j$下\Term{stable}{稳定}可知，只要$\epsilon(x)=n$，公式(1.22)就成立，而无论$i>n$时是否有$x^i=0$。

\SourcePage{41}
公式(1.22)有如下推广。设$K$为任意\Term{lambda-ring}{$\lambda$-环}，除此之外不作\Term{restriction}{限制}；对$x\in K$，置
\begin{equation}\gamma^n(x)=(-1)^n\sum_{i=0}^n(-1)^i\lambda^i(x+n)=\lambda^n(x+n-1)\tag{1.23}\end{equation}
（末两个表达式的相等由直接归纳得出）。特别地，$\gamma^0(x)=1$、$\gamma^1(x)=x$；再置
\begin{equation}\gamma_t(x)=\sum_{n\geq0}\gamma^n(x)t^n\in K[[t]],\tag{1.24}\end{equation}
由此得到公式
\begin{equation}\gamma_t(x)=\lambda_{t/(1-t)}(x),\tag{1.25}\end{equation}
它\Term{equivalence}{等价}于
\begin{equation}\lambda_s(x)=\gamma_{s/(1+s)}(x).\tag{1.25 bis}\end{equation}
这表明$\lambda^i$可用$\gamma^i$表达，并有$\gamma_t(x+y)=\gamma_t(x)\gamma_t(y)$，从而映射$\gamma^i$在$K$上定义一个新的\Term{lambda-ring}{$\lambda$-环}结构。公式(1.22)于是表明，当$K=\widetilde A$且$x\in\widetilde A$时，有
\begin{equation}\gamma^n(x-\epsilon(x))=[0,1+(-1)^{n-1}(n-1)!\,x^n+\cdots].\tag{1.22 bis}\end{equation}
因此，将$\widetilde A^n$中的元素$\gamma^n(x-\epsilon(x))$模$\widetilde A^{n+1}$\Term{reduced}{约化}，就能得到\Term{set-component}{分量}$x^n$，相差的因子为$(-1)^{n-1}(n-1)!$。

\textbf{\Term{chern-homomorphism}{陈同态}。}对满足上述条件的任意\Term{graded-ring}{分次环}$A$，我们要定义一个\Term{augmentation}{增广}环\Term{homomorphism}{同态}
\begin{equation}\operatorname{ch}:\widetilde A\longrightarrow\widehat{A\otimes\mathbb Q}.\tag{1.26}\end{equation}
此\Term{homomorphism}{同态}由以下条件唯一确定：它是加法\Term{homomorphism}{同态}，将1变为1，$\operatorname{ch}(x)$的正\Term{degree}{次数}\Term{set-component}{分量}由$x$的\Term{set-component}{分量}的有理系数等权\Term{universal-polynomial}{泛多项式}给出，并且
\SourcePage{42}
\begin{equation}\operatorname{ch}([1,1+x^1])=\exp x^1=\sum_{n\geq0}(x^1)^n/n!.\tag{1.27}\end{equation}
于是立即得到$\operatorname{ch}$的显式表达式：设$\eta$为$\widehat{A\otimes\mathbb Q}$的加法自\Term{homomorphism}{同态}，对\Term{degree}{次数}为$i$的$x^i$定义为
\begin{equation}\eta(x^i)=\frac{1}{(-1)^{i-1}(i-1)!}\,x^i;\tag{1.28}\end{equation}
则有
\begin{equation}\operatorname{ch}\left([n,1+\sum_{i\geq1}x^i]\right)
=n+\eta\left(\log\left(1+\sum_{i\geq1}x^i\right)\right).\tag{1.29}\end{equation}
由于这个公式可以“求逆”，至少当$A$仅有\Term{finite}{有限}个\Term{homogeneous-component}{齐次分量}时如此，所以在这一情形，$\operatorname{ch}$给出环$\widetilde A\otimes\mathbb Q$到环$\widehat{A\otimes\mathbb Q}$的\Term{isomorphism}{同构}（这完全阐明了$\widetilde A\otimes\mathbb Q$的结构）。

回顾Hirzebruch [1]：一个\Term{formal-power-series}{形式幂级数}$f\in\mathbb Q[[t]]$通过\Term{universal-polynomial}{泛多项式}公式定义一个加法\Term{homomorphism}{同态}$1+\widehat A^+\to1+(\widehat{A\otimes\mathbb Q})^+$，记为$\mathfrak T_f$；当$f(t)=t/(1-\exp(-t))$时，简记$\mathfrak T_f=\mathfrak T$（$\mathfrak T$为Todd的首字母）。将$\mathfrak T(x)$的定义扩充到$x\in\widetilde A$。由此得到下述结果。
\begin{SGAresult}{命题}{1.3}
设$N\in\widetilde A$形如$[q,1+N^1+\cdots+N^q]$（故$i>q$时$N^i=0$）。若置
\[\lambda_{-1}(N)=\sum_{i\geq0}(-1)^i\lambda^i(N)=\sum_{i=0}^{n}(-1)^i\lambda^i(N),\]
则有
\begin{equation}\operatorname{ch}(\lambda_{-1}(N))=(-1)^qN^q\mathfrak T(N^\vee)^{-1},\tag{1.30}\end{equation}
\SourcePage{43}
\Term{equivalence}{等价}地，
\begin{equation}\operatorname{ch}(\lambda_{-1}(N^\vee))=N^q\mathfrak T(N)^{-1}.\tag{1.30 bis}\end{equation}
\end{SGAresult}

\SGAsection{4}{运算$\lambda^p(N,x)$}
设$A$为以两列无限个\Term{indeterminate}{不定元}$\lambda^iN$、$\lambda^ix$（$i\geq1$）为变量、系数在$\mathbb Z$中的\Term{formal-power-series}{形式幂级数}环。置$\lambda^1N=N$、$\lambda^1x=x$，并在$A$上引入唯一的\Term{special-lambda-ring}{特殊$\lambda$-环}结构，使$\lambda^i(N)=\lambda^iN$、$\lambda^i(x)=\lambda^ix$，且运算$\lambda^i$“\Term{continuous}{连续}”（若仅考虑\Term{indeterminate}{不定元}$\lambda^iN,\lambda^ix$的\Term{polynomial}{多项式}，就得到“由$N$和$x$生成的自由\Term{lambda-ring}{$\lambda$-环}”；我们的环$A$是它的一个\Term{completion}{完备化}）。于是可构成
\begin{equation}\lambda_{-1}(N)=\sum_{i\geq0}(-1)^i\lambda^iN,\tag{1.31}\end{equation}
这是常数项为1的\Term{formal-power-series}{形式幂级数}，因而\Term{invertible}{可逆}；故可无歧义地用下式定义$A$中的元素$\lambda^p(N,x)$：
\begin{equation}\lambda^p(N,x)\lambda_{-1}(N)=\lambda^p(x\lambda_{-1}(N)).\tag{1.32}\end{equation}
\begin{SGAresult}{\Term{theorem}{定理}}{1.4}
设$J_q$为$A$中由$i>q$的各$\lambda^iN$生成的\Term{closed-ideal}{闭理想}；将$A/J_q$等同于各$\lambda^ix$及各$\lambda^jN$（$1\leq j\leq q$）的\Term{formal-power-series}{形式幂级数}环。则$\lambda^p(N,x)$模$J_q$\Term{reduced}{约化}后为一个\Term{polynomial}{多项式}，且关于变量$\lambda^ix$为权$p$的\Term{isobaric-polynomial}{等权多项式}。

设$N,F$为\Term{characteristic}{特征}0的\Term{algebraically-closed-field}{代数闭域}$k$上\Term{dimension}{维数}分别为$q,q'$的\Term{finite-dimensional-vector-space}{有限维向量空间}；以$N^{(p)}$表示$N^p$到$N$的映射$(a_1,\ldots,a_p)\mapsto a_1+\cdots+a_p$的核。
\SourcePage{44}
令\Term{symmetric-group}{对称群}$\mathfrak S_p$\Term{action}{作用}在$N^p$上，因而\Term{action}{作用}在$N^{(p)}$及$\bigwedge(N^{(p)})$上，再\Term{action}{作用}在$\bigotimes^pF$及$\bigwedge(N^{(p)})\otimes\bigotimes^pF$上。将这个\Term{vector-space}{向量空间}视为$\mathfrak S_p\times\operatorname{GL}(N)\times\operatorname{GL}(F)$上的\Term{graded-module}{分次模}，取它的“\Term{alternating-part}{交错部分}”，这是$\operatorname{GL}(N)\times\operatorname{GL}(F)=G$上的一个\Term{graded-module}{分次模}。于是有
\begin{equation}\lambda^p(N,F)=\underline E_G\left(\left(\bigwedge(N^{(p)})\otimes\bigotimes^pF\right)^{\mathrm{alt}}\right).\tag{1.33}\end{equation}
\end{SGAresult}
注意：若$M$为$G$-模的\Term{complex}{复形}，以$\underline E_G(M)$表示$M$各\Term{set-component}{分量}的类在$\underline K(G)$中的\Term{alternating-sum}{交错和}；这就是$M$的\Term{euler-poincare-characteristic}{欧拉—庞加莱示性数}。公式(1.33)左端$\lambda^p(N,F)$表示将变量$\lambda^iN,\lambda^iF$分别替换为$G$-模$\bigwedge^i(N),\bigwedge^i(F)$在$\underline K(G)$中的类后所得的环$\underline K(G)$中的元素。

\begin{proof}
用第二个公式(1.11)展开(1.32)的右端，可见要证明第一个断言，只需证明令$x=1$所得的相应断言；即证明在\Term{lambda-ring}{$\lambda$-环}$\mathbb Z[\lambda^1N,\ldots,\lambda^qN]$中（其中$i>q$时$\lambda^iN=0$），元素$\lambda^p(\lambda_{-1}(N))$可被$\lambda_{-1}(N)$整除。先验地，商
\[\lambda^p(N,1)=\lambda^p(\lambda_{-1}(N))/\lambda_{-1}(N)\]
是$1\leq i\leq q$的各$\lambda^iN$的\Term{formal-power-series}{形式幂级数}，位于这些变量的\Term{formal-power-series}{形式幂级数}环的\Term{field-of-fractions}{分式域}$K$中。另一方面，仍以$N$表示\Term{characteristic}{特征}0的\Term{algebraically-closed-field}{代数闭域}$k$上\Term{dimension}{维数}为$q$的\Term{vector-space}{向量空间}，考虑元素
\[\ell^p(N)=\underline E_G\bigl((\bigwedge(N^{(p)}))^{\mathrm{alt}}\bigr)\in\underline K(G),\]
此时$G=\operatorname{GL}(N)$。我们知道（第2节，\Term{theorem}{定理}1.3），$\underline K(G)$可等同于
\SourcePage{45}
$K$中由各$\lambda^iN$（$1\leq i\leq q$）及$\lambda^qN$的逆元生成的\Term{subring}{子环}。另一方面，在此环中容易验证
\[\ell^p(N)\lambda_{-1}(N)=\lambda^p(\lambda_{-1}(N)).\]
由于我们处于一个大域$K$中，且$\lambda_{-1}(N)\ne0$，可得$\ell^p(N)=\lambda^p(N,1)$；两端都属于$\underline K(G)$与各$\lambda^iN$的\Term{formal-power-series}{形式幂级数}环之交，因而都是各$\lambda^iN$的\Term{polynomial}{多项式}。

为在\Term{generic}{一般}情形证明(1.33)，验证右端在$\underline K(G)$中满足与(1.32)相应的公式，即它与$\lambda_{-1}(N)$之积为$\lambda^p(F\lambda_{-1}(N))$。同样，由$\underline K(G)$的特殊性可知，$\lambda^p(N,F)\in\underline K(G)$与$\lambda_{-1}(N)$之积也是$\lambda^p(F\lambda_{-1}(N))$。由于$\underline K(G)$为\Term{integral-ring}{整环}（参见第2节），且$\lambda_{-1}(N)\ne0$（参见同处），公式(1.33)就得到了证明。
\end{proof}
\textbf{注记。}a）当$F$的\Term{dimension}{维数}$q'\geq p$时，(1.33)刻画了关于$\lambda^iN$（$1\leq i\leq q$）及$\lambda^jx$（$1\leq j\leq p$）的\Term{polynomial}{多项式}$\lambda^p(N,x)$。

b）若域$k$不是\Term{algebraically-closed-field}{代数闭域}，(1.33)似乎仍成立\footnote{可用相同方法证明。}。若要用下面阐述的方法在非\Term{algebraically-closed}{代数闭}的\Term{characteristic}{特征}0域上证明黎曼—罗赫\Term{theorem}{定理}，就必须研究这个问题。此外，从$\mathfrak S_p$-模取其\Term{alternating-part}{交错部分}，只有在\Term{base-field}{基域}\Term{characteristic}{特征}为0时才有合理意义。

现在设$K$为任意\Term{special-lambda-ring}{特殊$\lambda$-环}，$N$为$K$中的元素，且$i$充分大时$\lambda^iN=0$。
\SourcePage{46}
于是，对每个$x\in K$及每个整数$p\geq1$，都可考虑元素$\lambda^p(N,x)$；它们\Term{representable}{可表}为各$\lambda^iN$与$\lambda^jx$的整系数\Term{universal-polynomial}{泛多项式}。有$\lambda^1(N,x)=x$，且公式(1.32)在$K$中成立。为表述运算$x\mapsto\lambda^p(N,x)$的性质，引入一个新的\Term{lambda-ring}{$\lambda$-环}$K_N$是方便的。作为\Term{commutative-ring}{交换环}，$K_N$由下述环添加\Term{identity-element}{单位元}得到：其底层\Term{additive-group}{加法群}为$K$，\Term{multiplication}{乘法}为
\begin{equation}x\mathbin{\cdot_N}y=xy\mu\qquad(\mu=\lambda_{-1}(N)).\tag{1.34}\end{equation}
从$K_N$到$K_N[[t]]$的映射$\lambda_t$定义为
\begin{equation}\lambda_t(n,x)=(1+t)^n\sum_{p\geq0}\lambda^p(N,x)t^p,\tag{1.35}\end{equation}
其中置$\lambda^0(N,x)=(1,0)$（$K_N$的\Term{identity-element}{单位元}）。我断言，配备这些运算后，$K_N$为\Term{special-lambda-ring}{特殊$\lambda$-环}。因此，需要对$X,Y\in K_N$验证公式
\[\lambda_t(X+Y)=\lambda_t(X)\lambda_t(Y),\qquad
\lambda_t(XY)=\lambda_t(X)\circ\lambda_t(Y),\qquad
\lambda_t(\lambda^i(X))=\lambda^i(\lambda_t(X)).\]
只需在$X,Y$属于$K$时验证这些公式；并且可看出，只需在$K$为由$X,Y,N$生成、且服从$i>q$时$\lambda^iN=0$这些关系的自由\Term{special-lambda-ring}{特殊$\lambda$-环}时验证即可。但将\Term{identity-element}{单位元}变为\Term{identity-element}{单位元}、且在$K$上\Term{restriction}{限制}为\Term{homomorphism}{同态}$x\mapsto x\mu$的群\Term{homomorphism}{同态}
\begin{equation}K_N\longrightarrow K\tag{1.36}\end{equation}
与\Term{ring-structure}{环结构}及映射$\lambda^i$相容，这可立即验证。在目前的情形，
\SourcePage{47}
它在$K$上的\Term{restriction}{限制}是\Term{injectivity}{单射}；由于$K$为\Term{special-lambda-ring}{特殊$\lambda$-环}，由此容易得出上述关系对$X,Y\in K$确实成立，从而完成了断言的证明。

根据公式(1.23)，对$x\in K$和$n\geq1$，自然置
\begin{equation}\gamma^n(N,x)=\sum_{i=0}^{n-1}\binom{n-1}{i}\lambda^i(N,x).\tag{1.37}\end{equation}
换言之，置$\gamma^n(N,x)=\gamma^n(x_N)$，其中$x_N$表示将$x$视为\Term{lambda-ring}{$\lambda$-环}$K_N$的元素。
\begin{SGAresult}{命题}{1.5}
设$A$为带\Term{unit}{单位}的交换\Term{graded-ring}{分次环}，$N$为$\widetilde A$中形如$N=[q,1+N^1+\cdots+N^q]$的元素，$x$为$\widetilde A$中的任意元素：$x=[\nu,1+\sum_{i\geq1}x^i]$。则对每个整数$j\geq0$，有$\gamma^{q+j}(N,x)\in\widetilde A^j$，且$\gamma^{q+j}(N,x)$的\Term{degree}{次数}$j$\Term{set-component}{分量}形如
\begin{equation}(\gamma^{q+j}(N,x))^{(j)}=(-1)^{j-1}(j-1)!\,
G_{q,j}(\nu,x^1,\ldots,x^j;N^1,\ldots,N^j),\tag{1.38}\end{equation}
其中$G_{q,j}$为整系数\Term{universal-polynomial}{泛多项式}，由下列条件刻画：
\begin{equation}G_{q,j}(\nu,x^1,\ldots,N^j)N^q
=(-1)^q(x*\lambda_{-1}(N))^{(q+j)}.\tag{1.39}\end{equation}
\end{SGAresult}
\begin{proof}
可假定$A$为\Term{indeterminate}{不定元}$x^i,N^j$（$1\leq j\leq q$）的整系数\Term{polynomial-ring}{多项式环}。结合(1.32)和(1.37)，显然有
\begin{equation}\gamma^n(N,x)\lambda_{-1}(N)=\gamma^n(x\lambda_{-1}(N)).\tag{1.40}\end{equation}
当$n=q+j$时，右端的\Term{filtration}{滤过}\Term{degree}{次数}至少为$q+j$（使用公式(1.22 bis)，并注意$x\lambda_{-1}(N)$的\Term{augmentation}{增广}为零），而$\lambda_{-1}(N)$的\Term{filtration}{滤过}\Term{degree}{次数}为$q$（公式(1.22)）；
\SourcePage{48}
因而$\gamma^n(N,x)$的\Term{filtration}{滤过}\Term{degree}{次数}为$j$，因为由(1.18)，此处与$\widetilde A$相伴的\Term{graded-ring}{分次环}为\Term{integral-ring}{整环}。使用前面的公式，立即得到$(\gamma^{q+j}(N,x))^{(j)}$的所述\Term{form}{形式}。
\end{proof}

\SGAchapter{RRR II}{凝聚代数层的类与陈类}{Classes de faisceaux algébriques cohérents et classes de Chern}
\SGAsection{1}{周理论}
本章始终固定一个\Term{base-field}{基域}$k$；为简单起见，假定它为\Term{algebraically-closed-field}{代数闭域}。不过，下面大部分内容，甚至可能全部内容，在没有这一\Term{restriction}{限制}时仍然成立。所考虑的\Term{algebraic-space}{代数空间}和\Term{regular-map}{正则映射}均定义在$k$上。

若一个\Term{algebraic-space}{代数空间}\Term{isomorphism}{同构}于某个\Term{projective-space}{射影空间}的\Term{locally-closed-subset}{局部闭子集}，就称它为\Term{quasi-projective}{拟射影的}。

设$X$为\Term{dimension}{维数}$n$的\Term{connected}{连通}、\Term{nonsingular-variety}{非奇异}、\Term{quasi-projective}{拟射影}簇。以$A(X)$表示并称为$X$的\Term{chow-ring}{周环}的是$X$上\Term{cycle}{闭链}模\Term{rational-equivalence}{有理等价}的类所成的\Term{graded-ring}{分次环}（$A^i(X)$由\Term{dimension}{维数}$n-i$的\Term{cycle}{闭链}的类构成）。$A(X)$确为\Term{graded-ring}{分次环}这一事实由周的下述\Term{theorem}{定理}得出：
\SourcePage{49}
\begin{SGAresult}{\Term{theorem}{定理}}{2.1}
$X$上的每个\Term{cycle}{闭链}都\Term{rational-equivalence}{有理等价}于一个\Term{nonsingular}{非奇异}\Term{cycle}{闭链}。给定$A^i(X)$与$A^j(X)$中的元素，可找到它们的\Term{nonsingular}{非奇异}代表$Y^{n-i}$与$Z^{n-j}$，使两者处处\Term{transverse-intersection}{横截相交}。
\end{SGAresult}
\footnote{由于作者的一处误解，这个\Term{theorem}{定理}所断言的内容超过了目前已知的结果。实际上，只使用周的“\Term{moving-lemma}{移动引理}”，它说可以在两个\Term{cycle}{闭链}各自的类中移动它们，使其交没有超额；参见Chevalley研讨班1958，\textit{Anneau de Chow et applications}。}
（若一个\Term{cycle}{闭链}的\Term{branch}{分支}都是\Term{nonsingular}{非奇异}\Term{subvariety}{子簇}，就称其为\Term{nonsingular}{非奇异}。若一个\Term{cycle}{闭链}的任一\Term{branch}{分支}与另一个\Term{cycle}{闭链}的任一\Term{branch}{分支}都处处\Term{transverse-intersection}{横截相交}，就称两个\Term{cycle}{闭链}处处\Term{transverse-intersection}{横截相交}；即两个相交的簇在交点处都\Term{nonsingular}{非奇异}，且它们在该处的切簇处于\Term{general-position}{一般位置}……）

设$f$为$X$到$Y$的\Term{morphism}{态射}（即\Term{regular-map}{正则映射}），其中$X,Y$均为\Term{quasi-projective}{拟射影}\Term{nonsingular-variety}{非奇异簇}；则$f$利用上述\Term{theorem}{定理}定义\Term{graded-ring}{分次环}\Term{homomorphism}{同态}
\begin{equation}f^*:A(Y)\longrightarrow A(X)\tag{2.1}\end{equation}
（\Term{cycle}{闭链}类的\Term{inverse-image}{逆像}）。这样，$A(X)$成为$X$的\Term{contravariant-functor}{反变函子}。

设$f$为$X^n$到$Y^m$的\Term{proper}{固有态射}（用Weil的术语，即$f$的图在$Y$上\Term{complete}{完备}，参见[3]）。于是$f$还定义加法\Term{homomorphism}{同态}
\begin{equation}f_*:A(X^n)\longrightarrow A(Y^m),\tag{2.2}\end{equation}
它保持\Term{cycle}{闭链}的\Term{dimension}{维数}，因而使\Term{degree}{次数}增加$m-n$。

\SourcePage{50}
因此，对\Term{proper}{固有}\Term{morphism}{态射}而言，忽略\Term{multiplication}{乘法}结构和\Term{graded}{分次}，$A(X)$为$X$的\Term{covariant-functor}{协变函子}。注意经典公式
\begin{equation}f_*(x\cdot f^*(y))=f_*(x)\cdot y.\tag{2.3}\end{equation}
我们还将使用周的\Term{chern-class}{陈类}理论。它为$X$上每个\Term{vector-bundle}{向量丛}$E$（仍假定$X$为\Term{nonsingular}{非奇异}\Term{quasi-projective}{拟射影}簇）配以\Term{chern-class}{陈类}$C^i(E)\in A^i(X)$（$i\geq1$），使其满足以下条件。置$C^0(E)=1$，并置
\begin{equation}\widetilde C(E)=\left[\operatorname{rank}E,\sum_{i\geq0}C^i(E)\right]\in\widetilde A(X)\tag{2.4}\end{equation}
（$\widetilde A(X)$的定义见第一章第3节）。\Term{chern-class}{陈类}理论所应满足的条件\Term{representable}{可表}为：
\begin{equation}\widetilde C(E)=\widetilde C(E')+\widetilde C(E'')\tag{2.5}\end{equation}
当\Term{vector-bundle}{向量丛}$E$是由$E''$\Term{extension}{扩张}$E'$得到的时成立（若置$C(E)=\sum_{i\geq0}C^i(E)$，此条件表为$C(E)=C(E')C(E'')$）。公式(2.5)使$\widetilde C$可扩充为从$\underline K(\mathscr S(X))$（参见第一章第2节例2）到$\widetilde A(X)$的加法\Term{homomorphism}{同态}：
\begin{equation}\widetilde C:\underline K(\mathscr S(X))\longrightarrow\widetilde A(X).\tag{2.6}\end{equation}
还要求$\widetilde C$为\Term{lambda-ring}{$\lambda$-环}的\Term{lambda-homomorphism}{$\lambda$-同态}；后一条件也写为
\begin{equation}\widetilde C(E\otimes F)=\widetilde C(E)\widetilde C(F),\qquad
\widetilde C(\bigwedge^iE)=\lambda^i\widetilde C(E)\tag{2.7}\end{equation}
（$E,F$为两个\Term{vector-bundle}{向量丛}；当然，$\widetilde C$自动与\Term{augmentation}{增广}相容）。设$D$为$X$上的\Term{divisor}{除子}，$L(D)$为其定义的\Term{vector-bundle}{向量丛}；要求
\SourcePage{51}
\begin{equation}C^1(L(D))=\ell(D),\qquad C^i(L(D))=0\quad(i\geq1)\tag{2.8}\end{equation}
（对任意\Term{cycle}{闭链}$Z$，以$\ell(Z)$表示它在$A(X)$中的类）。最后，要求各$C^i$具有\Term{functoriality}{函子性}，即若$f$为$X$到$Y$的\Term{morphism}{态射}，$E$为$Y$上的\Term{vector-bundle}{向量丛}，则
\begin{equation}C^i(f^!(E))=f^*(C^i(E)),\tag{2.9}\end{equation}
其中$f^!(E)$为$E$在$X$上的\Term{inverse-image}{逆像}。也可将(2.9)表述如下：丛的\Term{inverse-image}{逆像}概念定义\Term{augmented-lambda-ring}{增广$\lambda$-环}的\Term{lambda-homomorphism}{$\lambda$-同态}
\begin{equation}f^!:\underline K(Y)\longrightarrow\underline K(X)\tag{2.10}\end{equation}
\footnote{此处采用的记号$f^!$与\Term{duality-theory}{对偶理论}中的记号冲突（例如参见R. Hartshorne，\textit{Residues and Duality}，Lecture Notes in Math. 第20卷，Springer，1966）。因此本研讨班将它改记为$f^*$。}
（这里以$\underline K(X)$代替$\underline K(\mathscr S(X))$），而(2.9)表示下图交换：
\[
\begin{tikzcd}
K(Y)\arrow[r,"f^!"]\arrow[d,"\widetilde C_Y"'] & K(X)\arrow[d,"\widetilde C_X"]\\
A(Y)\arrow[r,"f^*"'] & A(X).
\end{tikzcd}
\]
（在\Term{signature}{符号}$\widetilde C$上以所在\Term{space}{空间}作下标。）
\begin{SGAresult}{\Term{theorem}{定理}}{2.2}存在满足条件(2.5)、(2.7)、(2.8)和(2.9)的\Term{chern-class}{陈类}理论。\end{SGAresult}
我们还承认下述命题。
\begin{SGAresult}{命题}{2.3}设$Y$为$X$的\Term{closed-subset}{闭子集}，$U$为其补集；若$x\in A(X)$在$U$上诱导零，则$x$有一个代表是\Term{support}{支撑}在$Y$上的\Term{cycle}{闭链}。\end{SGAresult}

\SourcePage{52}
\begin{SGAresult}{推论}{}
设$p=\dim X-\dim Y$。若$i<p$，则\Term{homomorphism}{同态}$A^i(X)\to A^i(U)$为\Term{bijection}{双射}。\Term{homomorphism}{同态}$A^p(X)\to A^p(U)$的核中每个元素都是$Y$的\Term{dimension}{维数}$n-p$的\Term{branch}{分支}的\Term{linear-combination}{线性组合}（所以若$Y$\Term{irreducible}{不可约}，则它是$\ell(Y)$的倍数）。
\end{SGAresult}

\SGAsection{2}{凝聚代数层的陈类的定义}
设$\mathfrak C$为\Term{abelian-category}{阿贝尔范畴}，$\mathfrak C_0$为$\mathfrak C$的一个子类。以$\underline K(\mathfrak C_0)$表示以$\mathfrak C_0$中元素的\Term{isomorphism}{同构}类为基的\Term{free-group}{自由群}，模去由形如$(E)-(E')-(E'')$的元素所\Term{generated-subgroup}{生成的子群}所得的\Term{quotient-group}{商群}，其中$E$是由$E''$\Term{extension}{扩张}$E'$得到的（$E,E',E''\in\mathfrak C_0$）。必须假定$\mathfrak C_0$的元素关于\Term{isomorphism}{同构}关系的类构成一个\Term{set}{集合}。以$\gamma_{\mathfrak C_0}(E)$表示$E\in\mathfrak C_0$在$\underline K(\mathfrak C_0)$中的类；无混淆之虞时省略下标$\mathfrak C_0$。群$\underline K(\mathfrak C_0)$是关于以下函数$E\mapsto f(E)$的\Term{universal-problem}{泛问题}的解：它定义在$\mathfrak C_0$上、取值于任意\Term{abelian-group}{阿贝尔群}$A$中，且为\Term{additive-function}{加性}，即每当$E$为$E'$与$E''$的\Term{extension}{扩张}时，$f(E)=f(E')+f(E'')$。

设$F$为$\mathfrak C$到$\mathfrak C'$的\Term{additive-functor}{加性函子}，$\mathfrak C'_0$为$\mathfrak C'$的子类，满足与$\mathfrak C_0$相同的条件；并假定对每个$A\in\mathfrak C_0$，各$R^qF(A)$都属于$\mathfrak C'_0$，且$q$充分大时为零（假定$\mathfrak C$中存在\Term{injective-resolution}{内射消解}，从而$F$的\Term{derived-functor}{导出函子}$R^qF$存在）。则$F$由下式定义$\underline K(\mathfrak C_0)$到$\underline K(\mathfrak C'_0)$的\Term{homomorphism}{同态}，仍记为$F$：
\begin{equation}F(\gamma_{\mathfrak C_0}(E))=\sum_{i\geq0}(-1)^i\gamma_{\mathfrak C'_0}(R^iF(E)).\tag{2.11}\end{equation}
更\Term{generic}{一般}地，若有$\mathfrak C$到$\mathfrak C'$的\Term{cohomological-functor}{上同调函子}$(F^i)$，
\SourcePage{53}
且对每个$E\in\mathfrak C_0$，各$F^i(E)$属于$\mathfrak C'_0$，除\Term{finite}{有限}个外都为零，就得到$\underline K(\mathfrak C_0)$到$\underline K(\mathfrak C'_0)$的一个\Term{homomorphism}{同态}。这些注记可推广到\Term{multifunctor}{多函子}；此外，“经典的”\Term{spectral-sequence}{谱序列}给出一些\Term{transitivity}{传递性}质，这里不予详述。

特别地，若$F$为\Term{exact-functor}{正合函子}，则公式(2.11)简化为
\begin{equation}F(\gamma_{\mathfrak C_0}(E))=\gamma_{\mathfrak C'_0}(F(E)).\tag{2.11 bis}\end{equation}
特别地，这适用于$\mathfrak C=\mathfrak C'$、$F$为恒等\Term{functor}{函子}、$\mathfrak C_0\subset\mathfrak C'_0$的情形：得到$\underline K(\mathfrak C_0)$到$\underline K(\mathfrak C'_0)$的\Term{canonical}{典范}\Term{homomorphism}{同态}。

\begin{SGAresult}{\Term{theorem}{定理}}{2.3}
设$\mathfrak C$为\Term{abelian-category}{阿贝尔范畴}，$\mathfrak C_0\subset\mathfrak C_1$为两个子类，使$\mathfrak C_1$中\Term{object}{对象}的\Term{isomorphism}{同构}类构成一个\Term{set}{集合}。假定满足以下条件：

（i）对$\mathfrak C$中每个\Term{exact-sequence}{正合列}$0\to A'\to A\to A''\to0$，若$A,A''$属于$\mathfrak C_0$（分别地，属于$\mathfrak C_1$），则$A'\in\mathfrak C_0$（分别地，$A'\in\mathfrak C_1$）。

（ii）设$u:A\to B$、$v:C\to B$，其中$A,B,C\in\mathfrak C_1$，$u$\Term{surjectivity}{满射}；则存在$L\in\mathfrak C_0$及\Term{homomorphism}{同态}$u':L\to C$、$v':L\to A$，使$u'$\Term{surjectivity}{满射}且$vu'=uv'$。

（iii）对每个$A\in\mathfrak C_1$，存在整数$d=d(A)$，使对任意由$\mathfrak C_0$的\Term{object}{对象}构成的$A$的\Term{left-resolution}{左消解}$L$，都有$Z_d(L)\in\mathfrak C_0$（$Z_d$表示\Term{dimension}{维数}$d$的\Term{cycle}{闭链}）。

在这些条件下，$\underline K(\mathfrak C_0)$到$\underline K(\mathfrak C_1)$的\Term{canonical}{典范}\Term{homomorphism}{同态}为\Term{isomorphism}{同构}。
\end{SGAresult}
\begin{proof}[证明原理]
设$A\in\mathfrak C_1$。由（ii）、（iii），存在由$\mathfrak C_0$中\Term{object}{对象}构成的$A$的\Term{finite}{有限}\Term{resolution}{消解}$L$；置
\SourcePage{54}
\[f(A)=\sum_{i\geq0}(-1)^i\gamma_{\mathfrak C_0}(L^i).\]
证明这样构成的$\underline K(\mathfrak C_0)$中的元素与所选\Term{resolution}{消解}$L$无关：若有第二个\Term{resolution}{消解}$L'$，可用第三个\Term{resolution}{消解}$L''$“\Term{covering}{覆盖}”$L,L'$，并配备$L''$到$L$及$L'$的满\Term{homomorphism}{同态}（由（i）、（ii）可做到这一点）；于是归结为证明$\underline E_{\mathfrak C_0}(L)=\underline E_{\mathfrak C_0}(L'')$。二者之差为$\underline E_{\mathfrak C_0}(N)$，其中$N$为满\Term{homomorphism}{同态}$L''\to L$的核；由（i），其各项属于$\mathfrak C_0$。而$N$为各\Term{dimension}{维数}均\Term{acyclic-complex}{无环的复形}，故$\underline E_{\mathfrak C_0}(N)=0$，所述结论遂得。（注意：对各项属于$\mathfrak C_0$的\Term{finite}{有限}\Term{complex}{复形}$L$，$\underline E_{\mathfrak C_0}(L)$表示其\Term{euler-poincare-characteristic}{欧拉—庞加莱示性数}$\sum_{i\geq0}(-1)^i\gamma_{\mathfrak C_0}(L^i)$。）

同样可证明，如此构成的函数$f:\mathfrak C_1\to\underline K(\mathfrak C_0)$是\Term{additive}{加性}的，故定义一个$\underline K(\mathfrak C_1)$到$\underline K(\mathfrak C_0)$的\Term{homomorphism}{同态}；显然，这个\Term{homomorphism}{同态}为$\underline K(\mathfrak C_0)$到$\underline K(\mathfrak C_1)$的\Term{canonical}{典范}\Term{homomorphism}{同态}之逆。
\end{proof}
\textbf{注记。}为验证（ii），只需验证$\mathfrak C_1$包含于$\mathfrak C$的一个子类$\mathfrak C_2$中，且后者满足两个条件：

（ii a）每个$A\in\mathfrak C_2$都\Term{isomorphism}{同构}于某个$L\in\mathfrak C_0$的商。

（ii b）若$u:A\to B$，且$A,B\in\mathfrak C_2$，则$\operatorname{Ker}u\in\mathfrak C_2$。

\begin{SGAresult}{推论}{}
设$X$为\Term{quasi-projective}{拟射影}\Term{algebraic-space}{代数空间}。设$\mathscr F(X)$为$X$上的\Term{coherent-algebraic-sheaf}{凝聚代数层}的\Term{category}{范畴}，$\mathscr S(X)$为\Term{locally-free}{局部自由}\Term{coherent-algebraic-sheaf}{凝聚代数层}的\Term{category}{范畴}（\Term{equivalence}{等价}于$X$上的\Term{vector-bundle}{向量丛}的\Term{category}{范畴}），$\mathscr F_0(X)$为满足以下条件的\Term{coherent-algebraic-sheaf}{凝聚代数层}$F$的\Term{category}{范畴}：对每个$x\in X$，$F_x$为\Term{finite}{有限}\Term{cohomological-dimension}{上同调维数}的$\mathscr O_x$-模
\SourcePage{55}
（因而有$\mathscr S(X)\subset\mathscr F_0(X)\subset\mathscr F(X)$）。则\Term{canonical}{典范}\Term{homomorphism}{同态}
\begin{equation}\underline K(\mathscr S(X))\longrightarrow\underline K(\mathscr F_0(X))\tag{2.12}\end{equation}
为\Term{isomorphism}{同构}。
\end{SGAresult}
我们验证\Term{theorem}{定理}2.3的条件。条件（i）与（iii）由\Term{cohomological-dimension}{上同调维数}的定义以及\Term{noetherian-local-ring}{诺特局部环}$\mathscr O_x$上\Term{finite-type}{有限型}\Term{projective-module}{投射模}为\Term{free-module}{自由模}这一事实得出。为验证（ii），证明$\mathfrak C=\mathscr F(X)$满足注记中的（ii a）、（ii b）。第二个条件是平凡的；因此只需证明下述结果。

\begin{SGAresult}{命题}{2.4（Serre）}
设$X$为\Term{quasi-projective}{拟射影}\Term{algebraic-space}{代数空间}。则$X$上每个\Term{coherent-algebraic-sheaf}{凝聚代数层}都\Term{isomorphism}{同构}于一个\Term{locally-free}{局部自由}\Term{algebra-sheaf}{代数层}的商。
\end{SGAresult}
若$X$为\Term{projective-space}{射影空间}的\Term{closed-subset}{闭子集}，则由[2]，此命题成立：因为$n$充分大时$F(n)$由其\Term{section}{截面}生成，这就是说$F$为$\mathscr O(-n)^k$的商。\Term{generic}{一般}情形，$X$为\Term{projective-space}{射影空间}某\Term{closed-subset}{闭子集}的\Term{open-subset}{开子集}，只需使用下面的结果（Cartier和Serre曾在特殊情形中得到它）。

\begin{SGAresult}{命题}{2.5}
设$F$为\Term{algebraic-space}{代数空间}$X$的\Term{open-subset}{开子集}$U$上的\Term{coherent-algebraic-sheaf}{凝聚代数层}。则$F$\Term{isomorphism}{同构}于$X$上某\Term{coherent-algebraic-sheaf}{凝聚代数层}的\Term{restriction}{限制}。
\end{SGAresult}
\begin{proof}
在包含$U$、且$F$可延拓于其上的各\Term{open-subset}{开集}上应用Zorn\Term{lemma}{引理}，可归结到$X$\Term{affine}{仿射}的情形。设$\overline F$为$X$上的\Term{algebra-sheaf}{代数层}，它在\Term{open-subset}{开集}$V$上的\Term{section}{截面}为$F$在$U\cap V$上的\Term{section}{截面}。
\SourcePage{56}
容易看出$\overline F$为\Term{quasi-coherent}{拟凝聚}的（即在每个\Term{affine}{仿射}\Term{open-subset}{开集}$V$上，因此在这里就是整个$X$上，由该\Term{open-subset}{开集}\Term{coordinate-ring}{坐标环}上的一个模定义，此模可以是\Term{finite-type}{有限型}，也可以不是）。若$U$本身\Term{affine}{仿射}，这是直接的，因为此时$\overline F$在\Term{open-subset}{开集}$V$上由视为$\Gamma(V,\mathscr O_X)$-模的$\Gamma(U\cap V,F)$定义。\Term{generic}{一般}情形，$U$为\Term{finite}{有限}个\Term{affine}{仿射}\Term{open-subset}{开集}$U_i$之并，故$\overline F$为以下显然的\Term{quasi-coherent}{拟凝聚}层\Term{homomorphism}{同态}的核：
\[\sum_i\overline{F|_{U_i}}\longrightarrow\sum_{i,j}\overline{F|_{U_i\cap U_j}},\]
因而自身为\Term{quasi-coherent}{拟凝聚}层。由于假定$X$\Term{affine}{仿射}，$\overline F$便由其\Term{section}{截面}生成。又因$\overline F|_U=F$\Term{coherent}{凝聚}，存在$\overline F$的\Term{finite}{有限}个\Term{section}{截面}，生成$\overline F|_U$。这些\Term{section}{截面}生成$\overline F$的一个\Term{coherent}{凝聚}\Term{subsheaf}{子层}，其在$U$上的\Term{restriction}{限制}\Term{isomorphism}{同构}于$F$。
\end{proof}
现在假定$X$为\Term{nonsingular}{非奇异}\Term{quasi-projective}{拟射影}簇。由\Term{syzygy-theorem}{合冲定理}，采用\Term{theorem}{定理}2.3推论的记号，有$\mathscr F_0(X)=\mathscr F(X)$。故在此情形有
\begin{equation}\underline K(\mathscr S(X))=\underline K(\mathscr F(X)).\tag{2.13}\end{equation}
这一\Term{isomorphism}{同构}可表述如下。
\begin{SGAresult}{\Term{theorem}{定理}2.3的推论}{2}
设$X$为\Term{nonsingular}{非奇异}\Term{quasi-projective}{拟射影}簇。对每个\Term{abelian-group}{阿贝尔群}$A$，$X$上\Term{vector-bundle}{向量丛}的取值于$A$的\Term{additive-function}{加性函数}$f$，与$X$上\Term{coherent-algebraic-sheaf}{凝聚代数层}的取值于$A$的\Term{additive-function}{加性函数}$g$之间，存在一一对应。与每个$f$对应的是唯一满足下式的$g$：
\SourcePage{57}
\begin{equation}f(E)=g(\mathscr O_X(E))\tag{2.14}\end{equation}
对$X$上每个\Term{vector-bundle}{向量丛}$E$成立，其中$\mathscr O_X(E)$表示$E$在$X$上的\Term{regular-section}{正则截面}的芽所成的层。
\end{SGAresult}
为计算任意\Term{coherent-algebraic-sheaf}{凝聚代数层}$F$的$g(F)$，取由形如$\mathscr O_X(E_i)$的层构成的$F$的\Term{finite}{有限}\Term{resolution}{消解}，于是有
\begin{equation}g(F)=\sum_{i\geq0}(-1)^if(E_i).\tag{2.14 bis}\end{equation}
特别地，若$X$\Term{connected}{连通}，第1节中曾考虑\Term{vector-bundle}{向量丛}$E$的一个\Term{additive-function}{加性函数}$\widetilde C(E)$，取值于群$\widetilde A=\widetilde A(X)$。由此得到\Term{coherent}{凝聚}层的一个\Term{additive-function}{加性函数}$\widetilde C(F)$；再取各\Term{set-component}{分量}，得到函数$F\mapsto C^i(F)\in A^i(X)$。$C^i(F)$称为\Term{coherent}{凝聚}层$F$的\Term{chern-class}{陈类}；原则上可用(2.14 bis)，由\Term{vector-bundle}{向量丛}的\Term{chern-class}{陈类}来计算它们（但此时右端应改用\Term{multiplication}{乘法}记号）。至于$\widetilde C(F)$沿$\mathbb Z$的\Term{set-component}{分量}，就是层$F$的\Term{rank}{秩}，它对应于\Term{integral-ring}{整环}上的模的\Term{rank}{秩}：若$X$\Term{affine}{仿射}（总可以归结到此情形），$F$由$X$的\Term{coordinate-ring}{坐标环}上的一个\Term{finite-type-module}{有限型模}定义，此时取这个模的\Term{rank}{秩}……
% 翻译块止：第二章第3节标题之前。
